% oct13-thm-moment-fluctuations: oct13.tex:865-886 (label thm:moment-fluctuations)
\paragraph{Paper excerpt.}
\begin{verbatim}
\begin{theorem}[Sum of the distances from the origin]\label{thm:moment-fluctuations}
Let $\gauge$ be the Euclidean norm, let $0<\drift<\nf1d$, and let $\Rmom(n)$ be as in~\eqref{eq:radial-moment-def}. Then the following hold.
\begin{enumerate}[label=\textup{(\roman*)}]
\item \underline{\emph{Central limit theorem}}: as $n\to\infty$,
\begin{equation*}
\frac{\sum_{x\in A_n}|x|-n/(2\drift)}{r_n^{\nf{(d+3)}{2}}}
\xrightarrow{\mathrm d}\mathcal N\Bigl(0,\frac{2d\omega_d}{\drift(d+2)(d+3)}\Bigr)
\end{equation*}
and
\begin{equation*}
\frac{\Rmom(n)-r_n}{r_n^{\nf{(3-d)}{2}}}
\xrightarrow{\mathrm d}\mathcal N\Bigl(0,\frac{2}{\drift d\omega_d(d+2)(d+3)}\Bigr)\, .
\end{equation*}
\item \underline{\emph{Law of the iterated logarithm}}: for each choice of sign, almost surely,
\begin{align}
\limsup_{n\to\infty}\frac{\pm\bigl(\sum_{x\in A_n}|x|-n/(2\drift)\bigr)}{r_n^{\nf{(d+3)}{2}}\sqrt{2\log\log n}}
&=\Bigl(\frac{2d\omega_d}{\drift(d+2)(d+3)}\Bigr)^{\nf12}\, ,\notag\\
\limsup_{n\to\infty}\frac{\pm(\Rmom(n)-r_n)}{r_n^{\nf{(3-d)}{2}}\sqrt{2\log\log n}}
&=\Bigl(\frac{2}{\drift d\omega_d(d+2)(d+3)}\Bigr)^{\nf12}\, .\label{eq:moment-radius-lil}
\end{align}
\end{enumerate}
\end{theorem}
\end{verbatim}
\paragraph{Lean statement.}
\begin{verbatim}
theorem CERW.LimitShapes.moment_fluctuations {d : ℕ} (hd : 2 ≤ d) :
    let ωd : ℝ := (volume (Metric.ball (0 : EuclideanSpace ℝ (Fin d)) 1)).toReal
    ∀ ε : ℝ, 0 < ε → ε < 1 / (d : ℝ) →
    let r : ℕ → ℝ := fun n => ((d + 1) * n / (2 * d * ε * ωd)) ^ ((1 : ℝ) / (d + 1))
    let v₁ : ℝ := 2 * d * ωd / (ε * (d + 2) * (d + 3))
    let v₂ : ℝ := 2 / (ε * d * ωd * (d + 2) * (d + 3))
    ∀ {Ω : Type u} [MeasurableSpace Ω] (μ : Measure Ω) [IsProbabilityMeasure μ]
      (X : ℕ → Ω → Site d), CERW.IsCERW μ ε X →
      let S : ℕ → Ω → ℝ := fun n ω =>
        (∑ x ∈ CERW.departureRange (X · ω) n, euclidNorm x - n / (2 * ε)) /
          r n ^ (((d : ℝ) + 3) / 2)
      let R : ℕ → Ω → ℝ := fun n ω =>
        (CERW.momentRadius (X · ω) n - r n) / r n ^ ((3 - (d : ℝ)) / 2)
      TendstoInDistribution S atTop id (fun _ => μ) (gaussianReal 0 v₁.toNNReal) ∧
      TendstoInDistribution R atTop id (fun _ => μ) (gaussianReal 0 v₂.toNNReal) ∧
      ∀ᵐ ω ∂μ, ∀ δ : ℝ, 0 < δ → ∀ σ : ℝ, σ = 1 ∨ σ = -1 →
        (∀ᶠ n : ℕ in atTop, σ * S n ω ≤ (Real.sqrt v₁ + δ) * Real.sqrt (2 * Real.log (Real.log
            n))) ∧
        (∃ᶠ n : ℕ in atTop, (Real.sqrt v₁ - δ) * Real.sqrt (2 * Real.log (Real.log n)) ≤ σ * S n
            ω) ∧
        (∀ᶠ n : ℕ in atTop, σ * R n ω ≤ (Real.sqrt v₂ + δ) * Real.sqrt (2 * Real.log (Real.log
            n))) ∧
        (∃ᶠ n : ℕ in atTop, (Real.sqrt v₂ - δ) * Real.sqrt (2 * Real.log (Real.log n)) ≤ σ * R n ω)
\end{verbatim}
\paragraph{Proof of the registered statement.}
The paper statement is textually unchanged from the earlier pinned source.
Apply the proved declaration \texttt{CERW.Frozen.moment\_fluctuations}
(node \texttt{thm-moment-fluctuations} of \texttt{ledger/manifest.yaml}) with the displayed
standing assumptions and parameters. Its entire conclusion and binder order
are identical to the registered statement, after changing only the exported
name. This supplies every asserted clause and preserves the scope of every
constant. The retained natural-language proof is
\texttt{ledger/nl/thm-moment-fluctuations.tex}; no new argument or premise is needed.
